\chapter{Cronograma de 24 meses}\label{cronograma-de-24-meses}

{
\begingroup
\small
\singlespacing
\setlength{\tabcolsep}{4pt}
\begin{longtable}[]{@{}
  >{\raggedright\arraybackslash}p{(\linewidth - 4\tabcolsep) * \real{0.12}}
  >{\raggedright\arraybackslash}p{(\linewidth - 4\tabcolsep) * \real{0.44}}
  >{\raggedright\arraybackslash}p{(\linewidth - 4\tabcolsep) * \real{0.44}}@{}}
\toprule\noalign{}
\begin{minipage}[b]{\linewidth}\raggedright
Meses
\end{minipage} & \begin{minipage}[b]{\linewidth}\raggedright
Atividade
\end{minipage} & \begin{minipage}[b]{\linewidth}\raggedright
Entrega e decisão
\end{minipage} \\ \addlinespace[4pt]
\midrule\noalign{}
\endhead
\bottomrule\noalign{}
\endlastfoot
1--2 & Revisão dirigida, atualização até a data de início e comparação
com trabalhos de 2026 & Matriz de originalidade; fixação definitiva da
pergunta. \\ \addlinespace[4pt]
3--5 & Reconstrução do TG, polarizações exatas e vínculos de
Fierz--Pauli & Caderno de derivação e benchmarks públicos. \\ \addlinespace[4pt]
6--8 & Resposta tensorial de PTA e ORFs & Implementação validada em
limites conhecidos. \\ \addlinespace[4pt]
9--11 & Simulações e comparação das três formas de análise & Primeiro
mapa de erros e custo computacional. \\ \addlinespace[4pt]
12--15 & Calibração, priori versus posterior, covariância e espectros &
Resultado principal sobre inferência de massa. \\ \addlinespace[4pt]
16--18 & Subconjunto tensor mais helicidade zero; testes de robustez &
Avaliação da identificação de polarizações. Se necessário, manter esse
resultado como extensão delimitada. \\ \addlinespace[4pt]
19--21 & Aplicação pública viável ou ampliação dos testes simulados;
redação do artigo & Manuscrito completo e código documentado. \\ \addlinespace[4pt]
22--24 & Submissão, dissertação e defesa & Artigo submetido, dissertação
e material reproduzível. \\
\end{longtable}
\endgroup
}

Pressupõe-se formação em relatividade, métodos matemáticos e
programação. A análise integral de resíduos de uma PTA se beneficia de
orientação ou colaboração com experiência em temporização. O estágio
inicial pode ser executado em uma estação de trabalho; campanhas
estatísticas devem ser dimensionadas após benchmarks e podem exigir
recursos institucionais.
